\SourceLeaf{407}%
wenn ich alle Materialien zu diesen Themen beisammenhabe. Das wird aber nicht so schnell gehen, teils wegen der Menge der Gegenstände, teils wegen der zunehmenden und immer wieder neu auftretenden Schwierigkeiten. Das Nachdenken darüber wird jetzt erneut durch einige Vorlesungen unterbrochen werden.
\par\Needspace{8\baselineskip}
\begin{center}\textsc{Gauss} an \textsc{Schumacher}. Göttingen, 21. November 1825.\end{center}
\indent\mbox{. . . . . . .} Seit einiger Zeit nehme ich einen Teil der allgemeinen Untersuchungen über die krummen Flächen wieder auf, die die Grundlage meines geplanten Werks über Höhere Geodäsie bilden sollen. Das Thema ist ebenso reichhaltig wie schwierig und lässt mich gegenwärtig gar nicht zu anderen Arbeiten kommen. Leider sehe ich, dass ich dabei \textls{sehr weit} ausholen muss, da auch das Bekannte in einer anderen, den neuen Untersuchungen angepassten Form entwickelt werden muss. Man muss dem Baum bis in alle seine Wurzelfasern folgen, und manches davon kostet mich wochenlanges angestrengtes Nachdenken. Vieles davon gehört sogar in die Geometria situs, ein noch fast ganz unbearbeitetes Feld. Der Wunsch, den ich bei meinen Arbeiten immer hatte, ihnen eine solche Vollendung zu geben, ut nihil amplius desiderari possit, erschwert sie mir natürlich außerordentlich, ebenso wie die Notwendigkeit, sie wegen anderer, damit nicht zusammenhängender Dinge oft zu unterbrechen. Wenn ich den Kopf voll davon habe, können Sie sich kaum vorstellen, wie anstrengend es für mich manchmal ist, mich vormittags nach einer schlaflosen Nacht, wie ich sie leider jetzt oft habe, mit frischem Kopf in das hineinzuversetzen, was ich meinen Zuhörern vorzutragen habe, und anschließend sofort wieder mit frischer Kraft zu meinen Überlegungen zurückzufinden. Mitunter werde ich aber durch einen glücklichen neuen Fund belohnt. So habe ich zum Beispiel die Verallgemeinerung des \textsc{Legendre}schen Theorems, nach dem auf der Kugel die Seiten proxime proportional zu den Sinuswerten der
